% A REGISTERED RESULT SECTION. No numeral outside \ccnum; see check_numbers.py.
% The story's figures are read from the release benchmark's run directories by
% the number generator (`lineup.story.*`), and its comparative
% sentences are held to them by the number check.
\section{A merge that reads well and is wrong}
\label{sec:problem}

\paragraph{Why this is a common problem.} Teams write the same knowledge down more than once. A support team keeps a pair of articles about the same failure: one written for every release, the other for a single customer whose incident looked the same. Nobody folds them together, because that is a careful, unrewarding job. Every sentence in each has to end up somewhere, or be dropped on purpose, by someone who can say why. This is the job people now hand to a language model, and the model is very good at the part a reader can see.

The part a reader cannot see is what was left out. A merge that drops a quarter of one source reads exactly like a merge that drops nothing: shorter, cleaner, no gap where the missing material used to be. Where the sources disagree, the value the model picked reads as decided, not guessed. A detail the model invented is written in the same confident voice as the real ones. None of these leaves a mark on the page. Finding them means reading the sources beside the output, which is the work the merge was supposed to save.

\paragraph{One real case.} A bike-share operator has a runbook for its field crew and another for its support desk. The first sends a jammed dock to the workshop after \ccnum{lineup.story.a_count} failed attempts, the second after \ccnum{lineup.story.b_count}. Asked to merge them, \ccnum{lineup.story.row_label} returned a runbook that reads well. It says \ccnum{lineup.story.a_count}, it never mentions \ccnum{lineup.story.b_count}, and nothing on the page shows that the sources disagreed. The model had not overlooked it. At the default setting it is told to choose one value where the sources disagree and to note the choice in a \emph{record}, a short note about a source sentence that it did not copy word for word. It wrote that record: ``\ccnum{lineup.story.reason}'' A reader of the runbook never sees it.

LLossless caught it (Figure~\ref{fig:story}). Its report lists the support desk's sentence, at line \ccnum{lineup.story.b_line} of its runbook, as contradicted by the merge, and ends the run with exit code \ccnum{lineup.story.exit}, which means a problem in the merged document. Which number is right is for a person to decide. Without the report, nobody would know there was a decision to make. Every pair in my tests comes with a \emph{reference merge}, a correct merge written in advance. The one for this pair keeps \ccnum{lineup.story.a_count} as well, because the field crew's runbook gives the reason: forcing a jammed dock again bends the pin. So the model chose well, and still nothing on its page tells a reader that a choice was made. Had it chosen \ccnum{lineup.story.b_count}, the page would have read just as well.

% Figure: the opening case. Placed in Section 1. Every value comes from the
% generated figures; check_numbers.py scans this file as a result figure.
\begin{figure}[H]
\centering
\begin{tikzpicture}[node distance=6mm and 8mm]
  \node[doc, text width=60mm, align=left] (a)
       {\textbf{Field crew's runbook}\\[2pt]Escalate \ccnum{lineup.story.code} to the workshop after
        \colorbox{goodfill}{\strut\textbf{\ccnum{lineup.story.a_count}}} failed attempts.};
  \node[doc, text width=60mm, align=left, right=8mm of a] (b)
       {\textbf{Support desk's runbook}\\[2pt]If the field crew have already tried
        \colorbox{badfill}{\strut\textbf{\ccnum{lineup.story.b_count}}} times, raise it to the
        workshop instead.};
  \node[model, below=9mm of $(a.south)!0.5!(b.south)$] (m)
       {\textbf{\ccnum{lineup.story.row_label}} merges them};
  \draw[flow] (a.south) -- (m.north west);
  \draw[flow] (b.south) -- (m.north east);
  \node[doc, text width=60mm, align=left, below=9mm of m, xshift=-34mm] (out)
       {\textbf{The merged runbook}\\[2pt]Escalate \ccnum{lineup.story.code} to the workshop after
        \colorbox{goodfill}{\strut\textbf{\ccnum{lineup.story.a_count}}} failed attempts.\\[2pt]
        \footnotesize\itshape No word of the other count.};
  \node[draw=docgrey, dashed, fill=white, rounded corners=2pt, text width=56mm,
        align=left, font=\footnotesize, inner sep=5pt, below=9mm of m, xshift=34mm] (rec)
       {\textbf{The record}, which no reader of the runbook sees\\[2pt]
        superseded: ``\ccnum{lineup.story.reason}''};
  \draw[flow] (m.south) -- (out.north);
  \draw[flow, dashed] (m.south) -- (rec.north);
  \node[prog, below=8mm of $(out.south)!0.5!(rec.south)$, text width=110mm] (rep)
       {\textbf{LLossless reports:} the support desk's sentence, line \ccnum{lineup.story.b_line},
        is contradicted by the merge. Exit code \ccnum{lineup.story.exit}: a person decides
        which count is right.};
  \draw[flow] (out.south) -- (rep.north -| out.south);
  \draw[flow] (rec.south) -- (rep.north -| rec.south);
\end{tikzpicture}
\caption{The opening case. The merge chose the field crew's count, which is the right
one, and noted the choice in a record that the merged runbook does not show. Had it
chosen the other count, the page would have read just as well. The report is what makes
the choice visible.}
\label{fig:story}
\floatfoot
\end{figure}


It was not one model's quirk. Of the \ccnum{lineup.story.rows} configurations in my benchmark (a configuration is one model, reached one way), \ccnum{lineup.story.dropped} merged the pair the same way. Whether a merge contradicts a source is a question about meaning, so LLossless puts it to a model, which I call the \emph{judge}. In the benchmark the judge was the model that had written the merge, and it reported the conflict on \ccnum{lineup.story.flagged} of those \ccnum{lineup.story.dropped} merges. On the others the model judging its own merge did not report it. That is why the tool keeps the checks that depend on a judge apart from the checks a program does alone (Section~\ref{sec:system}).

\paragraph{Declared loss and silent loss.} This paper depends on the difference between these, so I state it here. A \emph{declared loss} is content that the merge left out and said so, in a record. That is usually fine, and often what a merge should do. A reviewer can see it, and the tool checks the record. The runbook merge above is one: it left out the support desk's count and said so. A \emph{silent loss} is content that is missing with no record. Finding those is exactly what this tool is for. A program looks for every source sentence in the merge, and a sentence it cannot find and that no record explains is counted as a silent loss. That count is a place to look and not proof, because the program compares text and also flags a sentence that the merge kept in other words. By the count the runbook merge is clean, with \ccnum{lineup.story.silent} silent losses, because its loss was declared. LLossless failed the run all the same, because the judge found that the merge contradicts a source.

\paragraph{A model on guardrails.} The use of language models for real work is held back by trust and by cost. A model can state things that are not so, which is called hallucination, and it can quietly leave things out. And every use costs money. LLossless is built with both in mind.

On trust, the tool does not ask the model to be reliable. It puts the model on guardrails. The model does the writing. A program, which cannot hallucinate, checks everything that can be checked by comparing text: that every source sentence is accounted for, and that every number, link, file path and code block is unchanged. A model is asked only where meaning has to be judged. It must quote its evidence, and the tool checks that the quote is there. The run ends with an exit code, so an automated pipeline stops on a bad merge instead of publishing it. The limit should be stated just as plainly. The guardrails do not stop a model from making an error. They make the error visible and stop the run. Even that has limits. A program cannot tell a sentence in other words from a lost one, or a sentence that lost a clause from a reworded one, so a person has to look. And where a judge is needed, the judge can miss (Section~\ref{sec:system}).

On cost, every report has a line for what the run cost, or for why the cost is not known. Writing the merge and checking it are separate jobs in a run, and each job can go to a different model, so an expensive model can write while a cheaper one checks (Section~\ref{sec:usage}). A cheaper depth of checking exists, and Section~\ref{sec:depth} measures what it gives up. And the benchmark measured what a merge costs on each model. In it the configuration closest to the reference merges cost about a third of the most expensive one. It was also the cheapest of those that the benchmark did not disqualify (Section~\ref{sec:lineup}).

\paragraph{What the tool does.} The guardrails enforce one rule. Everything in the sources is either in the merge or declared as left out, and everything in the merge comes from a source. The tool merges documents with a model of the user's choice and then checks that rule in both directions. The checks that a program does alone, with no model, I call the \emph{mechanical checks}. They find a source sentence that is missing and has no record, and a changed number, link or piece of code. The checks that rest on the judge I call the \emph{judged checks}. They are made one statement at a time, with a quote for every verdict. The result is a report a person can check line by line and an exit code a pipeline can act on.

\paragraph{What I measured.} Because the mechanical checks are a program, they can score many models on the same documents in the same way. Section~\ref{sec:findings} reports one such benchmark, run in a single session that began on \ccnum{lineup.from}: \ccnum{lineup.rows} configurations of hosted models (models that their vendors run, reached over the network), each merging the same \ccnum{lineup.pairs} document pairs. The program counted a silent loss for \ccnum{lineup.flagged_rows} of the \ccnum{lineup.rows} on at least one pair. A count is only a place to look, so an AI model, and no person, read each counted sentence against its merge. On the first run of each pair it also read every sentence that the program had passed as reworded although no record names it. For \ccnum{lineup.disqualified} of the \ccnum{lineup.rows} the reading found content gone with no record, and the benchmark disqualifies those. One of them is out for a single sentence that lost its opening condition, a case the readings call borderline.

My first scoring rule let the count decide. The reading showed that most of what it had counted was still in the merge, so I amended the rule after the results were known. Both scorings are published, and Section~\ref{sec:audit} tells how the error was found. Neither scoring counts a loss that the merge declared, such as the dropped count in the runbook merge. Section~\ref{sec:limitations} comes back to that.

\paragraph{Contributions.}
\begin{points}
\item A tool that puts a language model on guardrails for one job, merging documents. The model writes the merge and the tool audits it. The mechanical checks ask no model, and they are kept apart from the judged checks, which need one.
\item A dated benchmark of current hosted models on merging without loss. It is published under the scoring as registered and under a scoring amended after the results, with the audit by reading that led from the first to the second.
\end{points}

\paragraph{Where the test documents come from.} The test documents are a handful of short texts written for this project, by me and by the coding agents, plus excerpts adapted from Wikipedia and from one NASA page. The repository says who wrote which and under what licence.

\noindent The terms this paper uses, with the meaning each has here, are listed in Appendix~\ref{sec:words}.
